% GENERATED FILE. Edit canonical lessons, then run npm run generate:books.
% canonical-source-hash: fnv1a64:cd7bf96a323cfc51
% canonical-lessons: ES-C755-calentamiento, ES-C755-trotar, ES-C755-atleta, ES-C755-lesionarse, ES-C755-musculo

\chapter{Calentar y trotar}
\label{ch:es-calentar-y-trotar}
\hlchaptermodality{755}

\begin{chapteropening}
\textbf{By the end of this chapter:} \emph{I can talk about warming up, jogging, muscles and sports injuries.}

Complete the last lesson of chapter 755: I can talk about warming up, jogging, muscles and sports injuries.
\end{chapteropening}

\section[el calentamiento]{el calentamiento — the warm-up; warming}

\label{lesson:ES-C755-calentamiento}



\begin{quote}
Before the new one: say the Spanish for gardening, then the Spanish for to sew.
\end{quote}

\subsection*{You'll want to know: el calentamiento}
\textbf{el calentamiento} — \textquotedblleft{}the warm-up; warming\textquotedblright{}.

\begin{quote}
Antes de correr hago diez minutos de calentamiento.
\end{quote}

\begin{quote}
\emph{Before running I do a ten-minute warm-up.}
\end{quote}

\emph{El calentamiento global} — global warming.

\begin{cousinweb}
\emph{calentamiento} is built on \emph{calentar} \textquotedblleft{}to warm\textquotedblright{}, made from \emph{caliente} \textquotedblleft{}hot\textquotedblright{} (EtymDB), which comes from Latin \emph{calēns}, from \emph{calēre} \textquotedblleft{}to be warm\textquotedblright{}.
\end{cousinweb}

\begin{grammarlens}[title={What you've built}]
One more word for telling, explaining and describing at B1.
\end{grammarlens}

\subsection*{Your turn}
\begin{itemize}
\raggedright
  \item \emph{Say it:} \emph{el calentamiento}
  \item \emph{Say it:} \emph{el calentamiento}, once more
  \item \emph{Say it:} a sentence of your own with \emph{el calentamiento}
  \item \emph{Recall:} say the Spanish for gardening, then the Spanish for to sew, then say \emph{el calentamiento} again
\end{itemize}

\begin{tcolorbox}[breakable,colback=teal!4,colframe=teal!35!black,title={Before you move on}]
What does el calentamiento mean? (\textquotedblleft{}The warm-up; warming\textquotedblright{}.) Say it once more.
\end{tcolorbox}
\section[trotar]{trotar — to jog, to trot}

\label{lesson:ES-C755-trotar}



\begin{quote}
Before the new one: say the Spanish for to sew, then the Spanish for the warm-up.
\end{quote}

\subsection*{You'll want to know: trotar}
\textbf{trotar} — \textquotedblleft{}to jog, to trot\textquotedblright{}.

\begin{quote}
Troto en el parque cada mañana.
\end{quote}

\begin{quote}
\emph{I jog in the park every morning.}
\end{quote}

Common in Latin America; in Spain people mostly say \emph{salir a correr}.

\begin{cousinweb}
\emph{trotar} was borrowed from French \emph{trotter} \textquotedblleft{}to trot\textquotedblright{}, which is usually traced to Frankish \emph{trottōn} \textquotedblleft{}to tread\textquotedblright{}.
\end{cousinweb}

\begin{grammarlens}[title={What you've built}]
One more word for telling, explaining and describing at B1.
\end{grammarlens}

\subsection*{Your turn}
\begin{itemize}
\raggedright
  \item \emph{Say it:} \emph{trotar}
  \item \emph{Say it:} \emph{trotar}, once more
  \item \emph{Say it:} a sentence of your own with \emph{trotar}
  \item \emph{Recall:} say the Spanish for to sew, then the Spanish for the warm-up, then say \emph{trotar} again
\end{itemize}

\begin{tcolorbox}[breakable,colback=teal!4,colframe=teal!35!black,title={Before you move on}]
What does trotar mean? (\textquotedblleft{}To jog, to trot\textquotedblright{}.) Say it once more.
\end{tcolorbox}
\section[el atleta]{el atleta — the athlete}

\label{lesson:ES-C755-atleta}



\begin{quote}
Before the new one: say the Spanish for the warm-up, then the Spanish for to jog.
\end{quote}

\subsection*{You'll want to know: el atleta}
\textbf{el atleta} — \textquotedblleft{}the athlete\textquotedblright{}.

\begin{quote}
Es una atleta muy rápida.
\end{quote}

\begin{quote}
\emph{She is a very fast athlete.}
\end{quote}

Same form for men and women: \emph{el atleta}, \emph{la atleta}.

\begin{cousinweb}
\emph{atleta} comes from Latin \emph{āthlēta}, from Greek \emph{athlēt\'{\={e}}s} \textquotedblleft{}one who competes for a prize\textquotedblright{} (EtymDB), from \emph{áthlos} \textquotedblleft{}contest\textquotedblright{}.
\end{cousinweb}

\begin{grammarlens}[title={What you've built}]
One more word for telling, explaining and describing at B1.
\end{grammarlens}

\subsection*{Your turn}
\begin{itemize}
\raggedright
  \item \emph{Say it:} \emph{el atleta}
  \item \emph{Say it:} \emph{el atleta}, once more
  \item \emph{Say it:} a sentence of your own with \emph{el atleta}
  \item \emph{Recall:} say the Spanish for the warm-up, then the Spanish for to jog, then say \emph{el atleta} again
\end{itemize}

\begin{tcolorbox}[breakable,colback=teal!4,colframe=teal!35!black,title={Before you move on}]
What does el atleta mean? (\textquotedblleft{}The athlete\textquotedblright{}.) Say it once more.
\end{tcolorbox}
\section[lesionarse]{lesionarse — to get injured}

\label{lesson:ES-C755-lesionarse}



\begin{quote}
Before the new one: say the Spanish for to jog, then the Spanish for the athlete.
\end{quote}

\subsection*{You'll want to know: lesionarse}
\textbf{lesionarse} — \textquotedblleft{}to get injured\textquotedblright{}.

\begin{quote}
Se lesionó la rodilla en el partido.
\end{quote}

\begin{quote}
\emph{He injured his knee in the match.}
\end{quote}

\emph{Una lesión} — an injury.

\begin{cousinweb}
\emph{lesionarse} is built on \emph{lesión} \textquotedblleft{}injury\textquotedblright{}, a learned borrowing from Latin \emph{laesiō} \textquotedblleft{}injury\textquotedblright{} (EtymDB), from \emph{laedere} \textquotedblleft{}to hurt\textquotedblright{}.
\end{cousinweb}

\begin{grammarlens}[title={What you've built}]
One more word for telling, explaining and describing at B1.
\end{grammarlens}

\subsection*{Your turn}
\begin{itemize}
\raggedright
  \item \emph{Say it:} \emph{lesionarse}
  \item \emph{Say it:} \emph{lesionarse}, once more
  \item \emph{Say it:} a sentence of your own with \emph{lesionarse}
  \item \emph{Recall:} say the Spanish for to jog, then the Spanish for the athlete, then say \emph{lesionarse} again
\end{itemize}

\begin{tcolorbox}[breakable,colback=teal!4,colframe=teal!35!black,title={Before you move on}]
What does lesionarse mean? (\textquotedblleft{}To get injured\textquotedblright{}.) Say it once more.
\end{tcolorbox}
\section[el músculo]{el músculo — the muscle}

\label{lesson:ES-C755-musculo}



\begin{quote}
Before the new one: say the Spanish for the athlete, then the Spanish for to get injured.
\end{quote}

\subsection*{You'll want to know: el músculo}
\textbf{el músculo} — \textquotedblleft{}the muscle\textquotedblright{}.

\begin{quote}
Me duelen los músculos después del gimnasio.
\end{quote}

\begin{quote}
\emph{My muscles ache after the gym.}
\end{quote}

\begin{cousinweb}
\emph{músculo} is a learned borrowing from Latin \emph{mūsculus} \textquotedblleft{}little mouse, muscle\textquotedblright{} (EtymDB), from \emph{mūs} \textquotedblleft{}mouse\textquotedblright{}; a moving muscle looked like a mouse under the skin.
\end{cousinweb}

\begin{grammarlens}[title={What you've built}]
One more word for telling, explaining and describing at B1.
\end{grammarlens}

\subsection*{Your turn}
\begin{itemize}
\raggedright
  \item \emph{Say it:} \emph{el músculo}
  \item \emph{Say it:} \emph{el músculo}, once more
  \item \emph{Say it:} a sentence of your own with \emph{el músculo}
  \item \emph{Recall:} say the Spanish for the athlete, then the Spanish for to get injured, then say \emph{el músculo} again
\end{itemize}

\begin{tcolorbox}[breakable,colback=teal!4,colframe=teal!35!black,title={Before you move on}]
What does el músculo mean? (\textquotedblleft{}The muscle\textquotedblright{}.) Say it once more.
\end{tcolorbox}
